\documentclass[11pt]{article}
\usepackage[margin=2.5cm]{geometry}
\usepackage{booktabs}
\usepackage{tabularx}
\usepackage{graphicx}
\usepackage[hidelinks]{hyperref}
\usepackage{times}
\usepackage{enumitem}
\usepackage{float}
\setlist[description]{leftmargin=2.2cm, labelwidth=2.0cm, itemsep=2pt, font=\normalfont\itshape}

\title{The Standard Subgraph Classification Datasets:\\ Four Real and Four Synthetic}
\author{}
\date{October 13, 2026}

\begin{document}
\maketitle

\section*{What these datasets are}

Subgraph classification has a standard test set. It came with SubGNN (Alsentzer et al.,
NeurIPS 2020), and almost every later method reports on it: GLASS (2022), S2N (2024), SSNP
(2023) and GPEN (2025). It has four real datasets. In each one there is a single large graph,
a list of subgraphs inside it, and one label per subgraph. SubGNN also made four synthetic
datasets: Density, Cut-Ratio, Coreness and Component. They have their own section below.

All four real datasets are already on this PC, in \texttt{data/raw/subgnn}. SubGNN did not
release the synthetic datasets, only code to make them. GLASS published the files it used,
and we downloaded those into \texttt{data/raw/subgnn\_synthetic}. Every number below was
measured from these files, unless it says it comes from a paper.

\begin{table}[H]
\centering
\small
\begin{tabularx}{\textwidth}{>{\raggedright\arraybackslash}p{2.4cm}>{\raggedright\arraybackslash}p{2.6cm}>{\raggedright\arraybackslash}X>{\raggedright\arraybackslash}X>{\raggedright\arraybackslash}p{2.1cm}}
\toprule
Dataset & Node & Subgraph (one color) & Label of the subgraph & Classes \\
\midrule
Pregnancy (ours) & a visit & a patient's visits & had the outcome & 2 \\
\addlinespace
PPI-BP & a protein & the proteins of one biological process & the kind of process & 6 \\
HPO-Metab & a symptom or a gene & the symptoms of one rare disease & the kind of metabolic disease & 6 \\
HPO-Neuro & a symptom or a gene & the symptoms of one rare disease & the kinds of neurological disease & 10, several each \\
EM-User & a workout & a user's workouts & the user's gender & 2 \\
\bottomrule
\end{tabularx}
\caption{How each dataset is split into subgraphs and what each subgraph is labelled with.}
\end{table}

\section*{Each dataset}

\begin{description}
\item[PPI-BP] Human proteins, with an edge when two proteins are known to interact. A
subgraph is the set of proteins that take part in one biological process, taken from the
Gene Ontology gene sets. The label is the broad group of that process: metabolism,
development, signal transduction, stress and death, cell organization, or transport.
\item[HPO-Metab] A knowledge graph of symptoms (from the Human Phenotype Ontology) and genes.
A subgraph is the list of symptoms of one rare disease. The label is one of six kinds of
metabolic disease, for example lipid or lysosomal. Each class has exactly 400 subgraphs.
\item[HPO-Neuro] The same graph as HPO-Metab, with neurological diseases instead. A disease
can be in more than one of the ten classes, so the label is a set of classes.
\item[EM-User] Workouts from the fitness app Endomondo. A subgraph is the workouts of one
user, and the label is the user's gender. It is the only binary dataset, and it is balanced:
$162$ women and $162$ men.
\end{description}

\section*{Size}

\begin{table}[H]
\centering
\small
\begin{tabular}{lrrrrrrr}
\toprule
 & Nodes & Edges & Avg.\ degree & Subgraphs & Avg.\ size & Median & Smallest--largest \\
\midrule
Pregnancy (ours) & $61{,}046$ & $972{,}557$ & $31.9$ & $6{,}590$ & $9.3$ & $5$ & $2$--$221$ \\
\addlinespace
PPI-BP & $21{,}521$ & $338{,}625$ & $31.5$ & $1{,}591$ & $10.2$ & $7$ & $1$--$160$ \\
HPO-Metab & $14{,}587$ & $3{,}238{,}174$ & $444.0$ & $2{,}400$ & $14.4$ & $13$ & $1$--$50$ \\
HPO-Neuro & $14{,}587$ & $3{,}238{,}174$ & $444.0$ & $4{,}000$ & $14.8$ & $14$ & $2$--$55$ \\
EM-User & $57{,}333$ & $4{,}573{,}417$ & $159.5$ & $324$ & $155.4$ & $134$ & $3$--$499$ \\
\bottomrule
\end{tabular}
\caption{Size of the graphs and of the subgraphs. Size is the number of nodes in a subgraph.
Edges are counted once, without self-loops.}
\end{table}

\section*{How the subgraphs sit in the graph}

\begin{table}[H]
\centering
\small
\setlength{\tabcolsep}{3pt}
\begin{tabular}{lrrrrrr}
\toprule
 & Of all nodes, & Of nodes in a & Pieces per & One piece & Edges inside & Same-label \\
 & in no subgraph & subgraph, in two+ & subgraph & & a subgraph & neighbours (chance) \\
\midrule
Pregnancy (ours) & $0\%$ & $0\%$ & $8.8$ of $9.3$ & -- & $0.4\%$ & $83.0\%$ ($81.7\%$) \\
\addlinespace
PPI-BP & $67.4\%$ & $47.8\%$ & $7.0$ of $10.2$ & $2.1\%$ & $2.3\%$ & $27.5\%$ ($19.4\%$) \\
HPO-Metab & $82.5\%$ & $78.2\%$ & $1.6$ of $14.4$ & $54.8\%$ & $2.9\%$ & $21.8\%$ ($17.0\%$) \\
HPO-Neuro & $76.3\%$ & $78.2\%$ & $1.5$ of $14.8$ & $59.3\%$ & $5.1\%$ & -- \\
EM-User & $87.6\%$ & $68.0\%$ & $30.6$ of $155.4$ & $0\%$ & $1.1\%$ & $49.6\%$ ($50.0\%$) \\
\bottomrule
\end{tabular}
\caption{The first two columns have different bases: the first is out of all nodes, the
second only out of the nodes that are in some subgraph. \emph{Pieces}: take the nodes of one subgraph, keep only the edges between them, and count
the separate groups that are left. A node with no edge to the rest of its subgraph is a
piece on its own. So ``$7.0$ of $10.2$'' means an average subgraph of $10.2$ nodes splits
into $7$ groups that do not touch each other. \emph{One piece} is the share of subgraphs that
do not split at all. \emph{Edges inside a subgraph} is the share of all edges in the graph
whose two ends are in the same subgraph. \emph{Same-label neighbours} looks at edges between two nodes of different subgraphs,
using only nodes that belong to one subgraph, and asks how often both subgraphs have the same
label; \emph{chance} is what random edges would give. It is left out for HPO-Neuro, where a
subgraph has several labels.}
\end{table}

\paragraph{Most nodes belong to no subgraph.} Between $67\%$ and $88\%$ of the nodes are only
background. In our data every visit belongs to a patient.

\paragraph{Subgraphs share nodes.} About half of the PPI-BP proteins that are in a subgraph
are in two or more, and about three quarters in HPO and EM-User. In our data a visit belongs
to exactly one patient.

\paragraph{PPI-BP and EM-User fall apart like ours.} A PPI-BP subgraph of $10.2$ proteins
splits into $7.0$ pieces, and only $2\%$ are in one piece. An EM-User subgraph splits into
$30.6$ pieces. The HPO subgraphs are mostly in one piece.

\paragraph{Neighbours carry some label signal, except in EM-User.} In PPI-BP, neighbouring
subgraphs share a label $27.5\%$ of the time, against $19.4\%$ by chance. In EM-User the two
numbers are the same, as in our data.

\paragraph{Size says little about the label.} In EM-User women and men have about the same
number of workouts ($152$ and $158$). In PPI-BP, signal transduction subgraphs are the
largest ($13.9$ proteins on average) and transport the smallest ($8.6$), so size helps a
little. In HPO-Metab lysosomal diseases are the largest ($18.2$ against $11.1$ for
carbohydrate).

\section*{The synthetic datasets}

Each one is a single made-up graph with subgraphs planted in it. The label of a subgraph
is a property of its shape, fixed when the dataset was made. The nodes have no features, so
a model can only use the structure.

\begin{description}
\item[Density] The label is the share of a subgraph's possible edges that are present:
low, medium or high. In the files the three ranges do not overlap: up to $0.074$, from
$0.079$ to $0.168$, and from $0.174$.
\item[Cut-Ratio] The label is how many edges leave the subgraph, out of how many could. Of
the edges that touch a subgraph, $77\%$, $80\%$ and $83\%$ leave it on average in the three
classes. The subgraphs themselves are almost complete: $95\%$ of their possible edges are
present.
\item[Coreness] The label is how deep in the core of the graph the subgraph sits, in three
classes. The average core number of its nodes is $19.4$, $25.1$ and $30.7$ in the three
classes, but each class covers almost the whole range, from about $3$ to $43$.
\item[Component] The label says whether the subgraph is one piece or several. All pieces,
in both classes, have exactly $15$ nodes. One class is a single piece of $15$ nodes. The
other is $5$ to $10$ pieces that do not touch each other, so $75$ to $150$ nodes.
\end{description}

\begin{table}[H]
\centering
\small
\begin{tabular}{lrrrrrrrr}
\toprule
 & Nodes & Edges & Avg.\ degree & Subgraphs & Classes & Avg.\ size & Median & Smallest--largest \\
\midrule
Density & $4{,}998$ & $29{,}962$ & $12.0$ & $250$ & 3 & $20.0$ & $20$ & $20$--$20$ \\
Cut-Ratio & $5{,}000$ & $84{,}295$ & $33.7$ & $250$ & 3 & $20.0$ & $20$ & $20$--$20$ \\
Coreness & $5{,}000$ & $119{,}205$ & $47.7$ & $221$ & 3 & $20.0$ & $20$ & $20$--$20$ \\
Component & $17{,}260$ & $91{,}619$ & $10.6$ & $250$ & 2 & $65.0$ & $15$ & $15$--$150$ \\
\bottomrule
\end{tabular}
\caption{Size of the synthetic graphs and of their subgraphs, counted as in Table 2.}
\end{table}

\begin{table}[H]
\centering
\small
\setlength{\tabcolsep}{3pt}
\begin{tabular}{lrrrrrr}
\toprule
 & Of all nodes, & Of nodes in a & Pieces per & One piece & Edges inside & Same-label \\
 & in no subgraph & subgraph, in two+ & subgraph & & a subgraph & neighbours (chance) \\
\midrule
Density & $55.0\%$ & $40.5\%$ & $8.1$ of $20.0$ & $3.2\%$ & $15.8\%$ & $32.9\%$ ($33.5\%$) \\
Cut-Ratio & $37.5\%$ & $42.1\%$ & $1.0$ of $20.0$ & $100\%$ & $53.4\%$ & $31.4\%$ ($33.5\%$) \\
Coreness & $11.6\%$ & $0\%$ & $1.0$ of $20.0$ & $100\%$ & $8.0\%$ & $38.6\%$ ($39.1\%$) \\
Component & $5.8\%$ & $0\%$ & $4.3$ of $65.0$ & $50.4\%$ & $93.4\%$ & $82.6\%$ ($84.1\%$) \\
\bottomrule
\end{tabular}
\caption{The same columns as Table 3. Here \emph{chance} is the sum of the squared label
shares among the ends of the edges counted.}
\end{table}

\paragraph{Size says nothing, or everything.} In Density, Cut-Ratio and Coreness every
subgraph has $20$ nodes. In Component size alone gives the label: $15$ nodes for one class,
$75$ or more for the other.

\paragraph{Neighbours carry no label signal.} In all four, neighbouring subgraphs share a
label as often as chance or a little less, as in our data.

\paragraph{Only Density falls apart like ours.} A Density subgraph of $20$ nodes splits into
$8.1$ pieces. Cut-Ratio and Coreness subgraphs are always one piece.

\paragraph{Coreness and Component do not share nodes.} No node is in two subgraphs, as in
our data. In Density and Cut-Ratio about $40\%$ of the nodes in a subgraph are in two or
more.

\paragraph{The test sets are small.} The GLASS code draws a new random split in every run:
half for training, a quarter for validation and a quarter for testing. That leaves $63$
test subgraphs ($55$ in Coreness), so one subgraph more or less right moves the score by
about $1.6$ points.

\section*{Published results}

\begin{table}[H]
\centering
\small
\begin{tabular}{lrrrr}
\toprule
 & PPI-BP & HPO-Metab & HPO-Neuro & EM-User \\
\midrule
SubGNN (2020) & $59.9 \pm 2.4$ & $53.7 \pm 2.3$ & $63.2 \pm 1.0$ & $81.4 \pm 4.6$ \\
GLASS (2022) & $61.9 \pm 0.7$ & $61.4 \pm 0.5$ & $68.5 \pm 0.5$ & $88.8 \pm 0.6$ \\
S2N, best version (2024) & $64.3 \pm 1.8$ & $63.9 \pm 1.7$ & $68.6 \pm 0.8$ & $89.0 \pm 1.6$ \\
\bottomrule
\end{tabular}
\caption{Micro-F1 on the test set, mean and spread over $10$ runs, from Table 2 of the S2N
paper (Kim and Oh, ICML 2024). For the binary EM-User this is the accuracy. The papers do not
report AUC, the metric we use.}
\end{table}

\paragraph{Scores differ a lot between datasets.} The best methods reach $89\%$ on EM-User
but only $64\%$ on PPI-BP and HPO-Metab, where there are six classes.

\begin{table}[H]
\centering
\small
\begin{tabular}{lrrrr}
\toprule
 & Density & Cut-Ratio & Coreness & Component \\
\midrule
SubGNN (2020) & $91.9 \pm 1.6$ & $62.9 \pm 3.9$ & $65.9 \pm 9.2$ & $95.8 \pm 9.8$ \\
GLASS (2022) & $93.0 \pm 0.9$ & $93.5 \pm 0.6$ & $84.0 \pm 0.9$ & $100.0 \pm 0.0$ \\
S2N+A (2024) & $93.2 \pm 2.6$ & $56.0 \pm 0.0$ & $85.7 \pm 5.8$ & $100.0 \pm 0.0$ \\
S2N+A with RWPE (2024) & $93.6 \pm 2.0$ & $89.2 \pm 2.6$ & $72.6 \pm 6.2$ & $100.0 \pm 0.0$ \\
\bottomrule
\end{tabular}
\caption{Micro-F1 on the synthetic test sets, mean and spread over $10$ runs, from Table 3
of the S2N paper. The SubGNN and GLASS rows are copied there from their own papers. S2N+A
is the GCNII version; RWPE adds random-walk position features to the nodes.}
\end{table}

\paragraph{Component is solved, Cut-Ratio is not.} Every method is at or near $100\%$ on
Component, where size gives the answer. Cut-Ratio spreads the methods most, from $56\%$ to
$93.5\%$.

\end{document}
